% Standalone version — no references to prior work framing
% All metrics extracted directly from result files — no values estimated or fabricated.
% UNSW-NB15: multiclass results now available (9 classes, Worms excluded, 1511 samples/class).
% UNSW-NB15 binary anon ablation and zero-day results still pending (notebooks not yet executed).

\section{Experimental Evaluation}\label{sec:evaluation}

%=======================================================================
\subsection{Dataset Overview}\label{subsec:datasets}
%=======================================================================

Table~\ref{tab:datasets} summarises the five IoT benchmark datasets used in
this evaluation. All datasets are subsampled to balanced class distributions
prior to training; the sample counts reflect the subsampled training corpus
size unless otherwise noted.

\begin{table*}[htbp]
\centering
\caption{Summary of IoT benchmark datasets used in evaluation.
  $^\dagger$Sample counts reflect balanced subsampled training corpora.
  $^\S$Only 370 normal flows exist in the Bot-IoT population; corpus is limited to 370 per class (740 total).
  $^\P$UNSW-NB15 multi-class corpus: 9 classes $\times$ 1{,}511 samples each (13{,}599 total);
  Worms class excluded (174 samples — insufficient for reliable per-class F1).
  A separate binary corpus of 10{,}000 samples is used for the anonymization ablation.}
\label{tab:datasets}
\begin{tabular}{llccl}
\toprule
\textbf{Dataset} & \textbf{Full Name} & \textbf{Classes} & \textbf{Samples}$^\dagger$ & \textbf{Representative Attack Families} \\
\midrule
CIC-IoT2023   & Canadian Institute for Cybersecurity IoT 2023 & 8  & 100{,}000 & DDoS-Flood, DoS-Flood, Mirai, MITM-ArpSpoofing, Recon, DNS\_Spoofing \\
WUSTL-IIoT    & WUSTL-IIoT-2021                               & 4  & 174{,}032 & Normal, Scanning, DoS, Reconn \\
TON\_IoT      & TON\_IoT Telemetry Dataset                    & 10 & 84{,}080   & Ransomware, Backdoor, DDoS, DoS, Injection, XSS, Password, Scanning, MitM, Normal \\
Bot-IoT       & Bot-IoT Dataset                               & 5  & 740$^\S$   & DDoS, DoS, OS\_Fingerprint, Service\_Scan, Keylogging \\
UNSW-NB15     & UNSW-NB15 Network Dataset$^\P$                & 9  & 13{,}599  & Normal, Generic, Exploits, Fuzzers, DoS, Reconnaissance, Analysis, Backdoor, Shellcode \\
\bottomrule
\end{tabular}
\end{table*}

%=======================================================================
\subsection{Multi-Class Classification Results}\label{subsec:multiclass_results}
%=======================================================================

Tables~\ref{tab:multiclass_llm_cic} and~\ref{tab:multiclass_llm_unsw} report
per-class and macro-averaged precision, recall, and F1 for the LLM-based
multi-class classifier on CIC-IoT2023 and UNSW-NB15 respectively.
Tables~\ref{tab:multiclass_rf_cic} and~\ref{tab:multiclass_rf_unsw} present
the corresponding Random Forest baselines. Binary classification results for
the remaining datasets are reported in Table~\ref{tab:binary_llm_summary}.

%-----------------------------------------------------------------
% CIC-IoT2023 — LLM multiclass
%-----------------------------------------------------------------
\begin{table*}[htbp]
\centering
\caption{LLM multi-class classification performance on CIC-IoT2023
  (100{,}000-sample balanced corpus, macro average across 8 classes).
  Best per-column value in \textbf{bold}.}
\label{tab:multiclass_llm_cic}
\begin{tabular}{lcccc}
\toprule
\textbf{Class} & \textbf{Precision} & \textbf{Recall} & \textbf{F1-Score} & \textbf{Support} \\
\midrule
BenignTraffic      & 0.2865 & 0.3426 & 0.3120 & 470 \\
DDoS-Flood         & 0.5163 & 0.4383 & 0.4741 & 470 \\
DDoS-Fragmentation & \textbf{0.7482} & \textbf{0.9400} & \textbf{0.8332} & 433 \\
DNS\_Spoofing      & 0.1065 & 0.2222 & 0.1440 &  81 \\
DoS-Flood          & 0.6370 & 0.1979 & 0.3019 & 470 \\
MITM-ArpSpoofing   & 0.3540 & 0.3053 & 0.3279 & 131 \\
Mirai              & 0.8068 & 0.9064 & \textbf{0.8537} & 470 \\
Recon              & 0.3024 & 0.4397 & 0.3584 & 141 \\
\midrule
\textbf{Macro Avg} & 0.4697 & 0.4740 & 0.4506 & 2666 \\
\textbf{Weighted Avg} & 0.5542 & 0.5300 & 0.5171 & 2666 \\
\textbf{Accuracy}  & \multicolumn{4}{c}{0.5300} \\
\bottomrule
\end{tabular}
\end{table*}

%-----------------------------------------------------------------
% CIC-IoT2023 — Random Forest multiclass baseline
%-----------------------------------------------------------------
\begin{table*}[htbp]
\centering
\caption{Random Forest multi-class baseline performance on CIC-IoT2023
  (100{,}000-sample corpus). Best per-column value in \textbf{bold}.}
\label{tab:multiclass_rf_cic}
\begin{tabular}{lcccc}
\toprule
\textbf{Class} & \textbf{Precision} & \textbf{Recall} & \textbf{F1-Score} & \textbf{Support} \\
\midrule
BenignTraffic      & 0.8577 & 0.9617 & 0.9067 & 470 \\
DDoS-Flood         & \textbf{1.0000} & 0.9957 & 0.9979 & 470 \\
DDoS-Fragmentation & 0.9931 & 0.9977 & 0.9954 & 433 \\
DNS\_Spoofing      & 0.8367 & 0.5062 & 0.6308 &  81 \\
DoS-Flood          & 0.9979 & \textbf{1.0000} & \textbf{0.9989} & 470 \\
MITM-ArpSpoofing   & 0.8684 & 0.7557 & 0.8082 & 131 \\
Mirai              & \textbf{1.0000} & 0.9957 & 0.9979 & 470 \\
Recon              & 0.7836 & 0.7447 & 0.7636 & 141 \\
\midrule
\textbf{Macro Avg} & 0.9172 & 0.8697 & 0.8874 & 2666 \\
\textbf{Weighted Avg} & 0.9505 & 0.9509 & 0.9487 & 2666 \\
\textbf{Accuracy}  & \multicolumn{4}{c}{0.9509} \\
\bottomrule
\end{tabular}
\end{table*}

%-----------------------------------------------------------------
% UNSW-NB15 — LLM multiclass
%-----------------------------------------------------------------

\begin{table*}[htbp]
\centering
\caption{LLM multi-class classification performance on UNSW-NB15
  (13{,}599-sample balanced corpus, 9 classes $\times$ 1{,}511 samples,
  Worms excluded). Best per-column value in \textbf{bold}.}
\label{tab:multiclass_llm_unsw}
\begin{tabular}{lcccc}
\toprule
\textbf{Class} & \textbf{Precision} & \textbf{Recall} & \textbf{F1-Score} & \textbf{Support} \\
\midrule
Analysis       & 0.0780 & 0.0364 & 0.0497 & 302 \\
Backdoor       & 0.1330 & 0.0894 & 0.1069 & 302 \\
DoS            & 0.0727 & 0.0265 & 0.0388 & 302 \\
Exploits       & 0.1103 & 0.0497 & 0.0685 & 302 \\
Fuzzers        & 0.5000 & 0.0232 & 0.0443 & 302 \\
Generic        & \textbf{1.0000} & 0.0199 & 0.0390 & 302 \\
Normal         & 0.0495 & \textbf{0.3013} & 0.0850 & 302 \\
Reconnaissance & \textbf{1.0000} & 0.0066 & 0.0132 & 302 \\
Shellcode      & 0.2180 & 0.1921 & \textbf{0.2042} & 302 \\
\midrule
\textbf{Macro Avg}    & 0.3513 & 0.0828 & 0.0722 & 2718 \\
\textbf{Weighted Avg} & 0.3513 & 0.0828 & 0.0722 & 2718 \\
\textbf{Accuracy}     & \multicolumn{4}{c}{0.0828} \\
\bottomrule
\end{tabular}
\end{table*}

%-----------------------------------------------------------------
% UNSW-NB15 — Random Forest multiclass baseline
%-----------------------------------------------------------------

\begin{table*}[htbp]
\centering
\caption{Random Forest multi-class baseline performance on UNSW-NB15
  (same corpus). Best per-column value in \textbf{bold}.}
\label{tab:multiclass_rf_unsw}
\begin{tabular}{lcccc}
\toprule
\textbf{Class} & \textbf{Precision} & \textbf{Recall} & \textbf{F1-Score} & \textbf{Support} \\
\midrule
Analysis       & 0.3135 & 0.3311 & 0.3221 & 302 \\
Backdoor       & 0.2625 & 0.2252 & 0.2424 & 302 \\
DoS            & 0.2635 & 0.2748 & 0.2690 & 302 \\
Exploits       & 0.5331 & 0.5596 & 0.5460 & 302 \\
Fuzzers        & 0.7173 & 0.7815 & 0.7480 & 302 \\
Generic        & \textbf{1.0000} & \textbf{0.9868} & \textbf{0.9933} & 302 \\
Normal         & 0.9365 & 0.7815 & 0.8520 & 302 \\
Reconnaissance & 0.8276 & 0.7947 & 0.8108 & 302 \\
Shellcode      & 0.8584 & \textbf{0.9636} & 0.9080 & 302 \\
\midrule
\textbf{Macro Avg}    & 0.6347 & 0.6332 & 0.6324 & 2718 \\
\textbf{Weighted Avg} & 0.6347 & 0.6332 & 0.6324 & 2718 \\
\textbf{Accuracy}     & \multicolumn{4}{c}{0.6332} \\
\bottomrule
\end{tabular}
\end{table*}

The UNSW-NB15 results expose a dataset-level limitation not present in CIC-IoT2023.
The LLM achieves a macro F1 of only $0.07$ versus Random Forest's $0.63$, a 9$\times$
gap compared to the 2$\times$ gap observed on CIC-IoT2023 ($0.45$ vs.\ $0.89$).
Crucially, the gap is not solely an LLM weakness: even Random Forest struggles on
Analysis~(0.32), Backdoor~(0.24), and DoS~(0.27), indicating that these classes
share overlapping traffic signatures at the feature level.
The root cause is convergence: multiple attack classes independently learn near-identical
volume-based rules (\texttt{sload}~$>$~$5\times10^7$, \texttt{rate}~$>$~$10^5$),
causing the proportion-of-rules-fired scorer to tie across classes and default to the
Normal label for the majority of test rows.
This manifests in the high precision but near-zero recall for Generic~(1.00/0.02) and
Reconnaissance~(1.00/0.007): the rules are correct when they fire, but fire so rarely
that $>$90\% of true instances are missed.
The result is interpretable as a diagnostic finding: UNSW-NB15's sub-class attack
families are statistically nearly indistinguishable in the feature space sampled,
and the zero-shot rule-generation approach cannot resolve this ambiguity without
gradient-based training on the discriminative directions between classes.

%-----------------------------------------------------------------
% Consolidated binary-mode LLM summary (WUSTL, TON, Bot, UNSW)
%-----------------------------------------------------------------

Binary classification results for WUSTL-IIoT, TON\_IoT, Bot-IoT, and UNSW-NB15
are reported in Table~\ref{tab:binary_llm_summary}. These serve as the basis
for the anonymization and zero-day analyses in subsequent sections.

\begin{table*}[htbp]
\centering
\caption{Binary LLM classification performance across datasets
  (macro-averaged F1, mean $\pm$ std across 3 seeds, 5 rule-refinement rounds).
  $^\S$Bot-IoT corpus limited to 740 samples (370 normal rows in full dataset); results should be interpreted with caution.}
\label{tab:binary_llm_summary}
\begin{tabular}{lccccc}
\toprule
\textbf{Dataset} & \textbf{$n$ samples} & \textbf{F1 (Original)} & \textbf{Std} & \textbf{F1 (Anonymized)} & \textbf{Std} \\
\midrule
CIC-IoT2023  & 100{,}000 & 0.7043 & 0.1089 & \textbf{0.7260} & 0.0938 \\
WUSTL-IIoT   & 174{,}032 & 0.8458 & 0.0778 & 0.6678 & 0.0513 \\
TON\_IoT     & 84{,}080  & 0.5770 & 0.1675 & \textbf{0.6743} & 0.0452 \\
Bot-IoT      & 740$^\S$  & 0.8736 & 0.0273 & \textbf{0.8936} & 0.0086 \\
UNSW-NB15    & 10{,}000 & 0.7779 & 0.0058 & 0.7683 & 0.0010 \\
\bottomrule
\end{tabular}
\end{table*}

%=======================================================================
\subsection{Zero-Day Detection Analysis}\label{subsec:zeroday_results}
%=======================================================================

Table~\ref{tab:zeroday_summary} reports zero-day detection performance for
each dataset. The ZDR column denotes the fraction of withheld-class flows
correctly flagged as anomalous; Known-F1 is the macro-averaged F1 on the
remaining known classes. A ZDR near 1.0 indicates that the LLM-generated rules
are discriminative enough to isolate the withheld attack family, while
Known-F1 measures the cost to closed-set performance of removing one class.

\begin{table*}[htbp]
\centering
\caption{Zero-day detection performance. The withheld class is excluded entirely
  from training and the retrieval corpus. ZDR~=~Zero-Day Detection Rate;
  Known-F1~=~macro F1 on remaining classes. All results at seed~42.
  $^\dagger$Two experimental runs available; the higher-ZDR run is reported;
  see Table~\ref{tab:zeroday_wustl_both} in the appendix for the comparison.}
\label{tab:zeroday_summary}
\begin{tabular}{llccc}
\toprule
\textbf{Dataset} & \textbf{Withheld Class} & \textbf{Zero-Day Samples} & \textbf{Known-F1} & \textbf{ZDR} \\
\midrule
CIC-IoT2023  & MITM-ArpSpoofing  & 11{,}211 & 0.8648 & 0.0218 \\
CIC-IoT2023  & Mirai-udpplain    & 31{,}864 & 0.8301 & 0.0004 \\
WUSTL-IIoT   & Reconn$^\dagger$  & 8{,}240  & \textbf{0.8005} & \textbf{0.9989} \\
TON\_IoT     & ransomware        & 14{,}735 & 0.6745 & \textbf{0.9999} \\
Bot-IoT      & Service\_Scan     & 58{,}626 & \textbf{0.9662} & 0.6490 \\
UNSW-NB15    & Reconnaissance    & 13{,}987 & 0.7547 & 0.4820 \\
\bottomrule
\end{tabular}
\end{table*}

The results span the full ZDR range. WUSTL-IIoT (Reconn) and TON\_IoT
(ransomware) achieve near-perfect ZDR ($\geq 0.999$) with only a modest
reduction in Known-F1, indicating that the withheld class occupies a distinct
region of feature space that the system can isolate without supervision. By
contrast, CIC-IoT2023 exhibits very low ZDR for both withheld classes
($\leq 0.022$), reflecting substantial overlap between withheld and known
traffic distributions. Bot-IoT achieves intermediate detection
(ZDR~$= 0.649$) at the cost of the highest Known-F1, suggesting that
Service\_Scan shares statistical signatures with known bot families.
UNSW-NB15 (Reconnaissance) yields a ZDR of $0.482$, the lowest of the
non-CIC scenarios, and is the only case where traditional ML classifiers
decisively outperform the LLM on zero-day detection
(DT: $0.845$, RF: $0.879$ versus LLM: $0.482$).
This reversal is consistent with the multiclass analysis: UNSW-NB15
Reconnaissance traffic shares volume-level statistical signatures with the
known attack classes in the training set, so the threshold rules learned from
known attacks fire on Reconnaissance flows at roughly chance rate.
The anonymization ablation for this experiment confirms feature names are
irrelevant: $\Delta$ZDR~$= +0.005 \pm 0.009$ across three seeds, indicating
that the low ZDR is structural rather than name-driven.

For WUSTL-IIoT, two runs with different rule configurations were available;
Table~\ref{tab:zeroday_wustl_both} in Appendix~\ref{app:detailed} presents
both results (ZDR of 1.0 and 0.999 respectively).

%=======================================================================
\subsection{Anonymization Ablation Findings}\label{subsec:anon_results}
%=======================================================================

Table~\ref{tab:anon_ablation} compares macro-averaged F1 under original
(semantic) versus anonymized (opaque-identifier) feature names across all
datasets with ablation results available. The Jaccard index quantifies
how consistently the system selects the same top-5 features across seeds when
feature names are hidden.

\begin{table*}[htbp]
\centering
\caption{Anonymization ablation results (macro-averaged F1, mean $\pm$ std
  across 3~seeds). $\Delta$F1~=~F1$_{\text{anon}}$~$-$~F1$_{\text{orig}}$.
  Jaccard: mean $\pm$ std feature-set overlap between original and anonymized
  conditions across seeds.}
\label{tab:anon_ablation}
\begin{tabular}{lcccccc}
\toprule
\textbf{Dataset} &
  \textbf{F1 (Original)} & \textbf{Std} &
  \textbf{F1 (Anonymized)} & \textbf{Std} &
  \textbf{$\Delta$F1} & \textbf{Jaccard} \\
\midrule
CIC-IoT2023 & 0.7043 & 0.1089 & \textbf{0.7260} & 0.0938 & $+$0.0217 & 0.889 $\pm$ 0.157 \\
WUSTL-IIoT  & \textbf{0.8458} & 0.0778 & 0.6678 & 0.0513 & $-$0.1780 & 0.083 $\pm$ 0.118 \\
TON\_IoT    & 0.5770 & 0.1675 & \textbf{0.6743} & 0.0452 & $+$0.0974 & 0.310 $\pm$ 0.084 \\
Bot-IoT     & 0.8736 & 0.0273 & \textbf{0.8936} & 0.0087 & $+$0.0201 & 0.587 $\pm$ 0.112 \\
UNSW-NB15   & 0.7779 & 0.0058 & 0.7683 & 0.0010 & $-$0.0096 & 0.508 $\pm$ 0.112 \\
\bottomrule
\end{tabular}
\end{table*}

Four of the five datasets show a $|\Delta\text{F1}| \leq 0.01$, demonstrating
that anonymization does not consistently degrade detection performance.
The exception is WUSTL-IIoT, where anonymization reduces F1 by $0.178$ and the
mean Jaccard drops to $0.083$, indicating that the system selects an almost
entirely different feature subset when semantic names are hidden. This suggests
that for WUSTL-IIoT the LLM relies on human-interpretable semantic cues to
locate discriminative flow-level features; without these cues it falls back on
network identity features that generalise less well.
UNSW-NB15 shows a small negative $\Delta\text{F1}$ of $-0.0096$ and a moderate
Jaccard of $0.508$, grouping it with CIC-IoT2023, TON\_IoT, and Bot-IoT as
datasets where anonymization has negligible impact.
Across these four datasets, Jaccard values of $0.508$--$0.889$ indicate moderate
to high feature-set consistency, confirming that the system discovers structurally
equivalent rules regardless of whether feature names are revealed.

%=======================================================================
\subsection{Threshold Sensitivity Analysis}\label{subsec:threshold}
%=======================================================================

Figures~\ref{fig:threshold_mitm}--\ref{fig:threshold_botiot} plot the
ZDR and FPR as a function of the SMCC detection threshold $\tau$ for each
zero-day scenario. Operating points that maximise the difference
ZDR~$-$~FPR are marked on each curve.

\begin{figure*}[htbp]
  \centering
  \includegraphics[width=0.48\textwidth]{../iot-llm-hids-pg/1-cic-iot/results/zeroday-MITM-ArpSpoofing-thresholds.png}
  \hfill
  \includegraphics[width=0.48\textwidth]{../iot-llm-hids-pg/1-cic-iot/results/zeroday-Mirai-udpplain-thresholds.png}
  \caption{Threshold sensitivity curves for CIC-IoT2023. Left: MITM-ArpSpoofing
    withheld (11{,}211 zero-day samples). Right: Mirai-udpplain withheld
    (31{,}864 zero-day samples). Both scenarios exhibit low ZDR at all
    thresholds, consistent with the high feature-space overlap between
    withheld and known classes.}
  \label{fig:threshold_mitm}
\end{figure*}

\begin{figure}[htbp]
  \centering
  \includegraphics[width=\columnwidth]{../iot-llm-hids-pg/2-wustl-iiot/results/zeroday-Reconn-thresholds.png}
  \caption{Threshold sensitivity for WUSTL-IIoT with Reconn withheld
    (8{,}240 zero-day samples). Near-perfect ZDR is maintained across a
    wide range of $\tau$, demonstrating that Reconn traffic is robustly
    separable from known classes in the generated rule space.}
  \label{fig:threshold_wustl}
\end{figure}

\begin{figure}[htbp]
  \centering
  \includegraphics[width=\columnwidth]{../iot-llm-hids-pg/3-ton-iot/results/zeroday-ransomware-thresholds.png}
  \caption{Threshold sensitivity for TON\_IoT with ransomware withheld
    (14{,}735 zero-day samples). ZDR reaches $\approx 1.0$ at moderate
    threshold values, confirming the discriminative power of the generated
    rule set for this class.}
  \label{fig:threshold_toniot}
\end{figure}

\begin{figure}[htbp]
  \centering
  \includegraphics[width=\columnwidth]{../iot-llm-hids-pg/4-bot-iot/results/zeroday-Service_Scan-thresholds.png}
  \caption{Threshold sensitivity for Bot-IoT with Service\_Scan withheld
    (58{,}626 zero-day samples). Intermediate ZDR ($0.649$) reflects partial
    overlap between Service\_Scan and known bot-traffic signatures.}
  \label{fig:threshold_botiot}
\end{figure}
